\documentclass[11pt]{amsart}
\usepackage{amsmath,amssymb,amsthm,mathtools}
\usepackage[margin=1in]{geometry}
\usepackage{enumerate}
\usepackage{hyperref}

\newtheorem{theorem}{Theorem}[section]
\newtheorem{proposition}[theorem]{Proposition}
\newtheorem{lemma}[theorem]{Lemma}
\newtheorem{corollary}[theorem]{Corollary}
\theoremstyle{definition}
\newtheorem{definition}[theorem]{Definition}
\newtheorem{problem}[theorem]{Problem}
\newtheorem{conjecture}[theorem]{Conjecture}
\theoremstyle{remark}
\newtheorem{remark}[theorem]{Remark}
\newtheorem{observation}[theorem]{Observation}

\newcommand{\RH}{\mathrm{RH}}
\newcommand{\bbR}{\mathbb{R}}
\newcommand{\bbC}{\mathbb{C}}
\newcommand{\bbZ}{\mathbb{Z}}
\newcommand{\bbQ}{\mathbb{Q}}
\DeclareMathOperator{\Real}{Re}
\DeclareMathOperator{\Imag}{Im}
\DeclareMathOperator{\Spec}{Spec}

\title[Finiteness and self-approximation]{Finiteness of off-line zeros and
sequential self-approximation of zeta:\\
a conditional reduction of the Riemann Hypothesis}

\author{David Alejandro Trejo Pizzo}
\thanks{Independent Researcher. \texttt{dtrejopizzo@gmail.com}}


\begin{document}

\begin{abstract}
Let $m\in\{0,1,2,\dots,\infty\}$ denote the number of quadruples of
non-trivial zeros of $\zeta$ off the critical line, so that the Riemann
Hypothesis (RH) is the statement $m=0$. We prove, by Bagchi's Rouch\'e
argument, a conditional dichotomy: if $\zeta$ satisfies sequential
self-approximation (Hypothesis~(SA): for every compact disk $D$ in
$\{\tfrac12<\Real s<1\}$ there are shifts $\tau_k\to\infty$ with
$\zeta(s+i\tau_k)\to\zeta(s)$ uniformly on $D$), then $m\in\{0,\infty\}$.
Consequently RH follows from Hypothesis~(SA) together with the finiteness
statement $m<\infty$ (Problem~(F)). Both hypotheses are open. Hypothesis~(SA)
follows from RH by Bagchi's theorem and is formally implied by Bagchi's
strong recurrence; it is not known whether it implies RH by itself.
Finiteness follows from RH and is not known to imply it. We also record an
elementary negative statement: for every $[\sigma_1,\sigma_2]\subset
(\tfrac12,1)$ and $p\ge1$, $t\mapsto\zeta(\cdot+it)$ is not Stepanov almost
periodic as an $L^p(\sigma_1,\sigma_2)$-valued function, a consequence of the
unboundedness of $\zeta$ on vertical lines. The remaining sections are a
heuristic discussion of the two hypotheses and of their relation to the
order ``finiteness before purity'' in Weil's and Deligne's proofs over
function fields. No proof of RH is claimed.
\end{abstract}

\maketitle
\tableofcontents
\newpage

\section{Introduction}

\subsection{A conditional splitting}
The heuristic synthesis \cite{P29} proposes that many approaches to RH meet
a common difficulty, the passage from averaged to individual statements
about zeros (the heuristic master quantifier of
\cite[Principle~3.1]{P29}). This paper isolates one elementary place where
RH splits into two statements of different kinds. Writing $m$ for the number
of off-critical zero quadruples,
\[
\RH \;=\; \underbrace{(m<\infty)}_{\text{finiteness}}
\;\wedge\;
\underbrace{(m<\infty\Rightarrow m=0)}_{\text{dichotomy}}.
\]
We prove that the dichotomy follows from a sequential self-approximation
hypothesis (Hypothesis~(SA), Definition~\ref{def:lp112}); the proof is
Bagchi's Rouch\'e step \cite{Bagchi87}, which under $m<\infty$ needs a single
replica instead of a positive-density family. Hypothesis~(SA) is implied by
RH and by strong recurrence. The finiteness statement $m<\infty$
(Problem~\ref{prob:L8}) is also implied by RH. We discuss, heuristically,
why finiteness looks like a statement of cohomological type, by analogy with
the function-field case, where $\dim H^1=2g$ is known before purity
\cite{Weil49,Del74}.

Both hypotheses are open, and neither is known to be equivalent to RH. The
contribution is the reduction (Theorem~\ref{thm:dichotomy},
Corollary~\ref{thm:architecture}) and a discussion of where each half
seems to be obstructed.

\subsection{Framework and notation}
Off-critical zeros of $\zeta$ come in quadruples
$\{\beta\pm i\gamma,\,1-\beta\pm i\gamma\}$ with $\beta=\tfrac12+\delta$,
$\delta\in(0,\tfrac12)$; $m$ counts the quadruples; $\RH\iff m=0$. By
Bohr--Landau, for every fixed $\sigma>\tfrac12$ the zeros with $\beta\ge\sigma$
have density zero among all zeros, $N(\sigma,T)=o(N(T))$
\cite[\S9.15]{Titchmarsh}; this does not exclude a positive proportion of
zeros with $\beta-\tfrac12\to0$. Let $Q$ be the Weil quadratic form on the
test space of \cite[Def.~2.1]{P25}. If
$0<m<\infty$, then $Q$ has negative index $\kappa(Q)=2m$ on the full test
space \cite[Prop.~4.7]{P25}; this is proved there under Hypothesis~D
(finitely many off-line quartets, \cite[Def.~2.4]{P25}).

\section{The dichotomy theorem}\label{sec:dichotomy}

\subsection{The self-approximation hypothesis}
\begin{definition}[Hypothesis~(SA): sequential self-approximation]\label{def:lp112}
For every compact disk $D\subset\{\tfrac12<\Real(s)<1\}$ there exists a
sequence $\tau_k\to\infty$ (which may depend on $D$) such that
$\zeta(s+i\tau_k)\to\zeta(s)$ uniformly on $D$.
\end{definition}

\begin{remark}[Quantifier form]
This is the per-disk form, and it is all the proof of
Theorem~\ref{thm:dichotomy} uses: a single disk around a single
off-critical zero. The form in which one sequence serves all disks
simultaneously is not needed and is not assumed.
\end{remark}

Hypothesis~(SA) demands a \emph{sequence} of good shifts, where Bagchi's
strong recurrence demands a set of good shifts of positive lower density.
Strong recurrence is equivalent to RH \cite{Bagchi87}
(see also \cite[\S5.5]{P27}); hence RH implies (SA), and strong recurrence
formally implies (SA). The converse direction is open
(\S\ref{ssec:notequiv}).

For disks free of zeros, (SA) holds unconditionally: if $\zeta\ne0$ on a
closed disk $D\subset\{\tfrac12<\Real s<1\}$, the universality theorem of Voronin
\cite{Voronin75}, in its form for compact sets with connected complement
\cite[Thm.~1.9]{Steuding}, applied to the non-vanishing function $\zeta|_D$, gives a
set of $\tau$ of positive lower density with
$\max_D|\zeta(s+i\tau)-\zeta(s)|<\varepsilon$. The content of (SA) lies in
the disks that contain zeros.

\subsection{The theorem}
\begin{theorem}[Conditional dichotomy]
\label{thm:dichotomy}
If Hypothesis~(SA) holds, then $m\in\{0,\infty\}$. Hence Hypothesis~(SA)
and $m<\infty$ together imply $\RH$.
\end{theorem}

\begin{proof}
Suppose $0<m<\infty$. Let $\gamma_{\max}$ be the largest $|\Imag\rho|$ over
the off-line zeros, and fix one of them, $\rho_0=\beta_0+i\gamma_0$ with
$\tfrac12<\beta_0<1$. Zeros are isolated, so we may choose
$0<r<\min(\beta_0-\tfrac12,1-\beta_0)$ with $\zeta\ne0$ on the circle
$|s-\rho_0|=r$. Let $D$ be the closed disk $|s-\rho_0|\le r$, which lies in
$\{\tfrac12<\Real s<1\}$, and put $\eta=\min_{\partial D}|\zeta|>0$. By
Hypothesis~(SA) there are $\tau_k\to\infty$ with
$\sup_D|\zeta(s+i\tau_k)-\zeta(s)|<\eta$ for $k\ge k_0$. By Rouch\'e's
theorem, $\zeta(\cdot+i\tau_k)$ has a zero in the interior of $D$ (as
$\zeta$ does, at $\rho_0$), so $\zeta$ has a zero $\rho'$ with
$\Real\rho'>\tfrac12$ and $\Imag\rho'\ge\gamma_0-r+\tau_k$. For $k$ large
this exceeds $\gamma_{\max}$, a contradiction. Hence $m=0$ or $m=\infty$.
\end{proof}

\begin{remark}[Why a single replica suffices]
Bagchi's argument \cite{Bagchi87} uses a positive-density family of
replicas. Under $m<\infty$ one replica above $\gamma_{\max}$ already
contradicts finiteness. Heuristically, this removes any conflict with the
zero-density estimates $N(\sigma,T)=o(T)$ \cite[Thm.~9.17]{Titchmarsh}, which do not obstruct
density-zero replication; what remains is that, under $\neg\RH$, the good
shifts would form an exceptional set (\S\ref{ssec:wall}).
\end{remark}

\subsection{Where the Euler product enters}\label{ssec:euler}
Finite $m>0$ occurs for functions with the symmetries of $\Xi$ but without an
Euler product. The Rouch\'e argument of Theorem~\ref{thm:dichotomy} applies
verbatim to such functions, so they fail the analogue of (SA); any proof of
(SA) for $\zeta$ must therefore use a property of $\zeta$ not shared by such functions, such as
its Dirichlet series or Euler product. For example, let
$0<b_j<\tfrac12$ and let $a_1<\dots<a_m$ be positive and distinct, put
$t_j=(j-\tfrac12)\pi$, and
\[
F(t)=\cos t\,\prod_{j=1}^m
\frac{(t^2-a_j^2)^2+2(t^2+a_j^2)b_j^2+b_j^4}{t^2-t_j^2}
=\cos t\,\prod_{j=1}^m\frac{(t^2-(a_j+ib_j)^2)(t^2-(a_j-ib_j)^2)}{t^2-t_j^2}.
\]
Since the $\pm t_j$ are simple zeros of $\cos t$, $F$ is entire, real on
$\bbR$, even, of order one, and its zeros are the real zeros of $\cos t$
other than $\pm t_j$ together with the $m$ quartets $\pm a_j\pm ib_j$. Thus
$G(s)=F(-i(s-\tfrac12))$ satisfies $G(s)=G(1-s)$, $\overline{G(\bar s)}=G(s)$,
and has exactly $m$ off-line quartets $\tfrac12\pm b_j\pm ia_j$, and the analogue of (SA) fails for $G$.

Heuristically, the Euler product of $\zeta$ supplies the ingredients of any
approach to (SA): the Besicovitch $B^2$ almost-periodicity of $\zeta(\sigma+it)$
for $\sigma>\tfrac12$ comes from the Dirichlet series; the $\bbQ$-linear
independence of $\{\log p\}$ makes the relevant almost-periods spread over
all heights; and the random Euler product model underlies the support
theorems behind universality. We do not make this precise.

\subsection{Hypothesis (SA) is not known to be equivalent to RH}\label{ssec:notequiv}
RH implies (SA). Conversely (SA) yields only the dichotomy: under $\neg\RH$
its consequence is $m=\infty$, which is not known to be contradictory. (SA)
is formally weaker than strong recurrence (sequence versus positive density,
a one-way implication). No material non-equivalence with RH can be proved
while RH is open, since under RH both statements are true; and nothing
known excludes that (SA) alone implies RH. Whether (SA) is easier to prove
than RH is an open question, not a claim of this paper.

\section{The finiteness problem}\label{sec:finiteness}

The discussion in this section is heuristic.

\subsection{No Siegel mechanism}
Could analysis prove $m<\infty$ directly, by an unconditional ``at most $C$
exceptions'' theorem as for Siegel zeros? Siegel-type repulsion works by
scarcity of a budget: the simple pole of $\zeta$ is a global $O(1)$ resource
near $s=1$, and two exceptional zeros compete for it. For off-critical
quadruples in the interior of the strip the local archimedean term is
$\asymp\log\gamma_0$, which grows, and there is no evident scarcity. In
positivity functionals of Weil type the contributions of distinct quadruples
add, so excluding two quadruples is not easier than excluding one. Finally,
statistics such as $N(T)$ and $S(T)$ do not distinguish a quadruple from a
double zero on the line at the same height (both contribute four zeros to the
count up to height $T$, two in the upper half-plane), so the information
carried by these averages does not see the abscissa.

\subsection{\texorpdfstring{Finiteness and $\dim H^1<\infty$}{Finiteness and dim H1 finite}}
Over function fields, the finiteness $\dim H^1=2g$ is available before
purity, and Deligne's amplification argument uses it \cite{Del74}; a
heuristic account of the amplification step and of what is missing for
$\zeta$ is in \cite[\S7]{P27}. The Riemann--Roch theorems of the
Connes--Consani programme (see \cite{CC21} for the archimedean positivity)
give, so far, values rather than signatures. In finite rank, finitely many
traces determine a finite set of eigenvalues by Newton's identities, whereas
\cite[Heuristic~4.3]{P29} describes why inertia is not determined by values
in general. This motivates stating the geometric half as a problem:

\begin{problem}[Problem~(F), finiteness]\label{prob:L8}
Prove $m<\infty$: the off-critical zeros of $\zeta$, if any, are finite in
number. By \cite[Prop.~4.7]{P25}, if this holds then the
Weil form $Q$ has finite negative index $\kappa(Q)=2m$ on the full test
space. Heuristically, one would expect a proof to come from a finiteness
statement of cohomological type over $\Spec\bbZ$ rather than from an
analytic estimate.
\end{problem}

\section{On sequential self-approximation}\label{sec:rigidity}

\subsection{What is known}
On $\Real s>1$, $\zeta(\sigma+it)$ is uniformly almost periodic in $t$
(Bohr \cite{Bohr22} studied almost-periodicity properties of Dirichlet series),
so self-approximation along a sequence is immediate there. For
$\tfrac12<\Real s<1$, the generalized strong recurrence $\mathrm{SR}_d$ (for every compact
$K\subset\{\tfrac12<\Real s<1\}$ with connected complement and every $\varepsilon>0$, the
set of $\tau$ with $\max_K|\zeta(s+id\tau)-\zeta(s+i\tau)|<\varepsilon$ has positive lower
density) is known for algebraic irrational $d$ (Nakamura \cite{Nak09}) and for all
irrational $d$ (Pa\'nkowski \cite{Pan10}); the case $d=0$ is Bagchi's strong
recurrence, equivalent to RH \cite{Bagchi87}; see \cite[\S3, \S5.5]{P27}. By
universality, (SA) holds for disks free of zeros (\S\ref{sec:dichotomy});
for disks containing a zero it is open. We are not aware of an unconditional
proof of sequential self-approximation of $\zeta$ on a disk containing a
zero in $\tfrac12<\Real s<1$.

\subsection{A negative statement in seminorm coordinates}\label{ssec:seminorm}
A natural attack on (SA) combines compactness (Montel) with
almost-periodicity ($B^2$). Heuristically, by mean-square bounds most shifts
keep $\zeta(\cdot+i\tau)$ bounded on a fixed disk $D$, and Montel's theorem
gives subsequential limits $g$; nothing identifies $g$ with $\zeta$. The
$B^2$ seminorm is a global average in $t$ and gives no information on a
fixed compact window. A seminorm that sees windows uniformly fails for a
simple reason:

\begin{theorem}[No Stepanov almost-periodicity, band sense]
\label{thm:stepanov}
For every $[\sigma_1,\sigma_2]\subset(\tfrac12,1)$, $\sigma_1<\sigma_2$, and
every $p\ge1$,
\[
\sup_{x\in\bbR}\int_x^{x+1}\!\!\int_{\sigma_1}^{\sigma_2}|\zeta(\sigma+it)|^p\,
d\sigma\,dt=\infty,
\]
so $t\mapsto\zeta(\cdot+it)$ is not $S^p$ almost periodic as an
$L^p(\sigma_1,\sigma_2)$-valued function. By the functional equation the same
holds for bands in $(0,\tfrac12)$.
\end{theorem}

\begin{proof}
Put $\sigma_0=(\sigma_1+\sigma_2)/2$ and
$r=\min((\sigma_2-\sigma_1)/2,\tfrac12)$. Since $\log|\zeta|$ is
subharmonic, so is $|\zeta|^p$, and for $s_0=\sigma_0+i(x+\tfrac12)$ the
disk $B(s_0,r)$ lies in $[\sigma_1,\sigma_2]\times[x,x+1]$, whence
\[
\int_x^{x+1}\!\!\int_{\sigma_1}^{\sigma_2}|\zeta|^p\,d\sigma\,dt
\ \ge\ \int_{B(s_0,r)}|\zeta|^p\ \ge\ \pi r^2\,|\zeta(s_0)|^p .
\]
The function $\zeta$ is unbounded on the line $\Real s=\sigma_0$
\cite[Thm.~8.12]{Titchmarsh}, so the right side is unbounded in $x$. An
$S^p$ almost periodic function is $S^p$-bounded, which proves the claim. For
bands in $(0,\tfrac12)$ the same argument applies, since
$\zeta(s)=\chi(s)\zeta(1-s)$ with $|\chi(\sigma+it)|\asymp|t|^{1/2-\sigma}\ge c>0$
for $0<\sigma<\tfrac12$, $|t|\ge1$, so $\zeta$ is unbounded on
$\Real s=\sigma_0\in(0,\tfrac12)$ as well.
\end{proof}

This is an elementary consequence of classical facts. The fixed-line scalar
version is not addressed here.

\begin{observation}[Seminorms, heuristically]\label{obs:seminorm}
Seminorms that see compact windows uniformly (uniform, Stepanov) are
infinite on $\zeta$ in the strip by the growth of $\zeta$; the seminorms
under which $\zeta$ is almost periodic ($B^2$) average over $t$ and assign
no weight to a fixed window. This is consistent with the heuristic master
quantifier of \cite[Principle~3.1]{P29}. We do not formulate or prove a
general statement of this kind.
\end{observation}

\subsection{The exceptional set of good shifts}\label{ssec:wall}
Under $\neg\RH$, let $\rho_0$ be an off-line zero and $D$ a small disk
around it as in the proof of Theorem~\ref{thm:dichotomy}. Each shift $\tau$
with $\sup_D|\zeta(s+i\tau)-\zeta(s)|<\eta$ produces, by Rouch\'e's theorem,
a zero $\rho'$ of $\zeta$ with $\Real\rho'\ge\sigma:=\beta_0-r>\tfrac12$ and
$|\Imag\rho'-\gamma_0-\tau|<r$, so each zero serves a set of shifts of
measure at most $2r$, and the set of good shifts in $[0,T]$ has measure at
most $2r\,N(\sigma,T+|\gamma_0|+r)$. By Carlson's zero-density estimate
$N(\sigma,T)=O(T^{4\sigma(1-\sigma)+\varepsilon})=o(T)$
\cite[Thm.~9.17]{Titchmarsh}, the set of good shifts has density zero.
Hypothesis~(SA) thus asks that an exceptional set be unbounded.

\begin{observation}[Cardinality and exceptionality, heuristically]\label{obs:cardinal}
Weakening a target from ``positive density of replicas'' to ``one replica''
removes the conflict with density theorems but, heuristically, not the
difficulty, because under $\neg\RH$ the witness lies in a density-zero set
that averaging methods do not reach. One could seek an intermediate
statement localising the mean-square closeness of a $B^2$ almost-period to a
fixed compact window; we do not formulate one.
\end{observation}

\section{The reduction}\label{sec:architecture}

\begin{corollary}[Two-hypothesis reduction]\label{thm:architecture}
If $m<\infty$ (Problem~\ref{prob:L8}) and Hypothesis~(SA)
(Definition~\ref{def:lp112}) holds, then $\RH$ holds. Both hypotheses are
open. Hypothesis~(SA) follows from RH (by Bagchi's theorem) and is formally
implied by strong recurrence; it is not known whether it implies RH by
itself. The bound $m<\infty$ follows from RH; it is not known to imply RH.
\end{corollary}

\begin{proof}
This is Theorem~\ref{thm:dichotomy}.
\end{proof}

\begin{remark}[The function-field order, heuristically]
Over function fields, finiteness ($\dim H^1=2g$) is established before
purity \cite{Weil49,Del74}, and given finiteness, finitely many traces
(Lefschetz numbers) determine the characteristic polynomial. In the
reduction above, given $m<\infty$ the replication argument of
Theorem~\ref{thm:dichotomy} needs only one witness. This analogy is only
suggestive.
\end{remark}

\begin{remark}[A contrast with the de~Bruijn--Newman constant]
Rodgers and Tao proved $\Lambda\ge0$ \cite{RT20} using statistics of all
zeros, which are visible to averages. The parameter $m$ is of a different
kind: off-line zeros have density zero in every fixed half-strip
$\Real s\ge\sigma>\tfrac12$, so they are invisible to such averages. This
heuristic contrast is what motivates splitting RH at $m<\infty$.
\end{remark}

\section{Summary}\label{sec:assessment}

What is proved here: Theorem~\ref{thm:dichotomy} (an elementary Rouch\'e
argument, Bagchi's replication step), Corollary~\ref{thm:architecture}, and
Theorem~\ref{thm:stepanov} (an elementary consequence of the unboundedness of
$\zeta$ on vertical lines in the strip), together with the finite-$m$
example of \S\ref{ssec:euler} and the remark that (SA) holds on zero-free disks by
universality. Everything else (Section~\ref{sec:finiteness}, the opening
paragraph of \S\ref{ssec:seminorm}, \S\ref{ssec:wall},
Observations~\ref{obs:seminorm} and~\ref{obs:cardinal}, the remarks on the
Euler product) is heuristic discussion.

What is not proved: either hypothesis. Problem~(F) and Hypothesis~(SA) are
open, each is implied by RH, and neither is known to be equivalent to RH.

\begin{thebibliography}{99}

\bibitem{Bagchi87} B.~Bagchi, \emph{Recurrence in topological dynamics and
the Riemann hypothesis}, Acta Math. Hungar. \textbf{50} (1987), 227--240.

\bibitem{Bohr22} H.~Bohr, \emph{\"Uber eine quasi-periodische Eigenschaft
Dirichletscher Reihen mit Anwendung auf die Dirichletschen $L$-Funktionen},
Math. Ann. \textbf{85} (1922), 115--122.

\bibitem{CC21} A.~Connes and C.~Consani, \emph{Weil positivity and trace
formula, the archimedean place}, Selecta Math.\ (N.S.) \textbf{27} (2021),
no.~77.

\bibitem{Del74} P.~Deligne, \emph{La conjecture de Weil. I},
Publ. Math. IH\'ES \textbf{43} (1974), 273--307.

\bibitem{Nak09} T.~Nakamura, \emph{The joint universality and the
generalized strong recurrence for Dirichlet $L$-functions}, Acta Arith.
\textbf{138} (2009), 357--362.

\bibitem{Pan10} \L.~Pa\'nkowski, \emph{Some remarks on the generalized
strong recurrence for $L$-functions}, in: New Directions in
Value-Distribution Theory of Zeta and $L$-Functions, Shaker Verlag, Aachen,
2009, 305--315.

\bibitem{RT20} B.~Rodgers, T.~Tao, \emph{The de Bruijn--Newman constant is
non-negative}, Forum Math. Pi \textbf{8} (2020), e6.

\bibitem{Steuding} J.~Steuding, \emph{Value-Distribution of $L$-Functions},
Lecture Notes in Math. \textbf{1877}, Springer, 2007.

\bibitem{Titchmarsh} E.~C.~Titchmarsh, \emph{The Theory of the Riemann
Zeta-Function}, 2nd ed., revised by D.~R.~Heath-Brown, Oxford Univ. Press,
1986.

\bibitem{Voronin75} S.~M.~Voronin, \emph{Theorem on the ``universality'' of
the Riemann zeta-function}, Izv. Akad. Nauk SSSR Ser. Mat. \textbf{39}
(1975), 475--486.

\bibitem{Weil49} A.~Weil, \emph{Numbers of solutions of equations in finite
fields}, Bull.\ Amer.\ Math.\ Soc.\ \textbf{55} (1949), 497--508.

\bibitem{P25} D.~A.~Trejo Pizzo, \emph{The Weil form and the scaling
operator: a Pontryagin-index dictionary in zero coordinates}, Paper~25 of
this series.

\bibitem{P27} D.~A.~Trejo Pizzo, \emph{One-dimensional propagation and the
Riemann Hypothesis: a heuristic survey of obstructions and two open
directions}, Paper~27 of this series.

\bibitem{P29} D.~A.~Trejo Pizzo, \emph{One wall in many coordinates: a
heuristic synthesis of the propagation trichotomy and a master quantifier
for the Riemann Hypothesis}, Paper~29 of this series.

\end{thebibliography}

\end{document}
